\section{관련 연구}\label{sec:related}
\subsection{팔레트 인지·위치 추정과 취급 응용}
팔레트 인지 연구에서 사용되는 입력은 컬러 영상, 깊이 영상, 점군 및 거리 센서 등으로 다양하다. RGB-D 기반 팔레트 연구는 깊이 정보를 이용한 템플릿 기반 인식과 위치 추정을 다룬다\cite{xiao2017}. 3차원 카메라 기반 연구는 3차원 카메라를 이용한 위치 추정을 다루며\cite{molter2018}, 창고 팔레트의 구조와 템플릿을 활용하는 RGB-D 방법\cite{zhao2022}, 능동 양안 점군의 특징과 정합을 이용하는 방법\cite{shao2023}, 깊이 카메라의 검출 영역과 점군 처리를 결합하는 방법\cite{hu2025}도 제안되었다. 이 연구들은 깊이 관측을 이용한다는 점에서 현재의 단안 영상 보정과 사용 정보가 다르다.

학습 기반 팔레트 위치 추정에서도 출력의 의미는 동일하지 않다. RGB 영상과 깊이 데이터를 결합한 검출·위치 추정\cite{li2021} 및 여러 깊이 기반 인지 모델의 비교\cite{rocha2023}는 센서 처리 경로 전체를 다룬다. 다중 작업 학습 연구는 영상에서 각도와 거리의 범주를 함께 예측하는 다중 작업 구성을 제안한다\cite{mok2021}. 구성요소 분할을 통한 위치 추정\cite{wu2025}과 위성항법 및 비전을 결합한 위치·회전 추정\cite{meng2024}은 각각 물체 표현과 좌표 정보의 다른 선택을 보여준다. 범주 분류 정확도, 거리 오차 및 전체 6D 자세 오차를 하나의 성능 순위로 해석하지 않는 이유가 여기에 있다.

팔레트의 인지 결과를 실제 취급으로 연결하는 연구에는 야외 팔레트 handling\cite{iinuma2020}, RGB-D 기반 포크 삽입\cite{iinuma2021d}, 다중 카메라를 이용한 적층\cite{iinuma2021s}, 팔레트 트럭의 검출·도킹\cite{tsiogas2021}, 학습 기반 화물 이송 시스템\cite{ren2022}, 영상 측정 기반 접근·삽입\cite{kita2026}이 있다. 본 연구는 이러한 전체 제어 절차를 새로 제안하는 대신, 제어에 공급되는 인지 결과의 보정에 초점을 둔다. 실제 장비에서 수집된 영상을 이용한다는 사실과 제안 보정기로 장비를 제어했다는 주장은 구분한다.

\subsection{합성 학습과 단안 팔레트 자세 추정}
Euro 팔레트 합성 학습 연구는 합성 영상을 이용해 DOPE 계열 추정기를 학습하고 RGB 카메라 기반 팔레트 자세 추정을 평가하였다\cite{knitt2022}. 팔레트 합성 위치 추정 연구는 합성 데이터와 팔레트 면의 기하 정보를 사용하는 접근이다\cite{mueller2024}. 따라서 합성 영상으로 학습하거나 코너와 알려진 기하를 이용해 자세를 구하는 전체 구성 자체는 선행연구와 구별되는 발명이 아니다.

관측 외형과 기하 표현을 분리하는 연구도 있다. 합성 학습과 stereo depth의 표면 기하를 이용하는 팔레트 인지\cite{beleznai2023}, disparity 및 점·선분 투표를 이용하는 연구\cite{beleznai2025}는 가림과 외형 변화를 다룬다. 앞면 기반 자세 추정은 앞면 관측과 신경망·커널 회귀를 이용해 적재물 외형 변화에 대응한다\cite{kai2025}. 단안 거리 기반 자세 추정은 단안으로 예측한 거리 정보와 기하 단서를 활용한다\cite{miura2026}. 물리적 깊이 센서의 관측, 단안으로 학습한 depth, 사전에 제공한 치수는 서로 다른 정보이므로 이를 명시한 비교가 필요하다.

\subsection{구조 정보를 이용한 키포인트 추정}
팔레트의 규칙적인 구조를 이용하는 접근은 이미 연구되어 왔다. 전이학습 기반 키포인트 연구는 개선된 YOLOv8s-pose 계열과 E형 단면의 키포인트, 강건한 기하 계산을 결합한다\cite{zhou2025jsc}. 개선된 인체 자세 모델의 팔레트 적용은 개선된 YOLOv11s-pose와 topology 정보를 결합한 위치 추정이다\cite{zhou2025nyas}. 하이퍼그래프 기반 팔레트 자세 추정은 키포인트의 고차 공간 관계와 구조 사전정보를 활용한다\cite{ye2026}. 따라서 본 연구의 동기는 기존 연구가 구조를 전혀 사용하지 않는다는 주장에 있지 않다. 이미 학습된 추정기의 결과를 유지하면서, 별도의 국소 보정 단계에서 치수 정보와 대응 규칙을 활용했을 때의 추가 효과를 확인하는 데 있다.


\begin{table*}[t]
\centering\footnotesize
\caption{주요 팔레트 선행과 본 연구의 입력·수정 대상. 센서와 초기화가 다른 보고 수치를 직접 순위화하지 않는다.}\label{tab:literature}
\begin{tabularx}{\textwidth}{p{.23\textwidth}p{.23\textwidth}Y}
\toprule 연구 & 핵심 정보 & 접근과 본 연구의 차이 \\\midrule
Euro pallet 합성 학습\cite{knitt2022}; Synthetic localisation\cite{mueller2024} & 합성 RGB, 코너와 객체 기하 & 초기 검출·자세 추정을 학습; 본 연구는 얻어진 초기 코너 이후의 추가 수정을 평가 \\
Front Face Shot\cite{kai2025}; Geometric cues\cite{beleznai2025} & 앞면 표현 또는 양안 시차·점·선분 & 관측 표현·센서 증거를 바꾸는 접근과 후단 코너 보정을 구분 \\
개선 pose 추정기\cite{zhou2025jsc,zhou2025nyas}; Hypergraph\cite{ye2026} & 단면 키포인트, 구조·고차 관계 & 초기 특징·키포인트 추정기를 개선; 물체 구조 활용 자체는 이미 알려진 원리 \\
Occlusion-Robust\cite{vu2024}; Monocular Metric Depth\cite{miura2026} & RGB-D 또는 단안 추정 깊이 & 관측 깊이·추정 깊이·등록 치수는 다른 정보 \\
Multi-Stage Domain-Adapted\cite{elmoaqet2025} & 영상 적응, 초기 자세, 렌더링 기반 보정 & 실제 자료 보정 평가의 참조 자세에 잡음을 더한 초기화와, 본 연구의 추정기 실제 출력은 다름 \\
제안 방법 & 초기 코너, 내부 영상 특징, 등록 치수 & 제한된 2D 코너 이동을 학습하고 같은 자세 계산으로 전후 효과를 평가 \\
\bottomrule
\end{tabularx}
\end{table*}


\subsection{입력 적응, 추정기 업데이트 및 출력 보정}
입력 표현을 바꾸는 접근의 예로 CUT는 대응 영상 쌍 없이 패치 수준의 대조 학습을 이용한 영상 변환을 제안한다\cite{cut2020}. 팔레트 연구에서는 CUT와 CosyPose를 결합한 다단계 보정이 보고되었다\cite{elmoaqet2025}. 이 방법은 영상 적응과 자세 보정을 함께 사용하므로, 합성 학습만으로 목표 영역의 정보를 전혀 보지 않은 설정과 동일하지 않다. 자기지도 pose 적응의 예인 Self6D도 단안 추론과 실사 RGB-D 학습 신호를 구분하여 이해해야 한다\cite{self6d2020}. 이러한 선행의 성공이 현재의 의사 레이블만을 사용하는 자기학습 설정에서도 같은 효과를 보장하지는 않는다.

출력 보정에서는 PoseFix가 이미지와 초기 키포인트를 입력받아 정제된 자세를 산출한다\cite{posefix2019}. DeepIM은 초기 자세로 렌더링한 영상과 관측 영상을 정합하여 자세를 반복적으로 갱신하며\cite{deepim2018}, CosyPose는 단일 시점 자세 가설과 객체 대칭을 다루고 다중 시점 장면을 정제한다\cite{cosypose2020}. 본 연구는 이들과 달리 기존 추정기의 내부 특징을 읽고 코너 주변의 제한된 변위 분포를 계산한다. 이는 전체 RGB 특징을 별도로 추출하거나 물체를 반복 렌더링하는 구성과 비교할 수 있지만, 내부 특징에 접근해야 하므로 모든 추정기에 같은 가중치를 그대로 적용할 수 있다고 주장하지 않는다. 특히 PoseFix의 원래 설정은 다른 추정기의 내부 코드나 특징을 필요로 하지 않는 반면, 본 방법은 기반별 특징·좌표 연결부와 보정 헤드 학습을 요구한다. 따라서 본문의 검증 대상은 동일 보정 절차의 여러 추정기 적용성이다.

분포에서 좌표를 계산하는 기대값 연산은 Integral Human Pose Regression 등의 선행에서 다루어졌다\cite{integral2018}. BB8은 코너 투영으로 자세를 표현하는 과정에서 회전 대칭이 학습에 일으키는 모호성을 다룬다\cite{bb82017}. 본 연구는 기대값 연산이나 대칭 고려 자체를 새롭게 제안하는 것이 아니라, 팔레트 코너 보정의 정보·감독·평가 조건을 묶어 추가 효과를 검토한다. 표~\ref{tab:literature}는 비교에 필요한 정보와 초기화 조건을 정리한다. 깊이 센서·렌더링·독립 초기 자세의 사용 여부가 다르면 보고된 정확도 수치를 현재 실험의 직접 기준선으로 옮기지 않는다.
